\section*{Bab \ref{ch:b1:crystals-ionic-covalent} --- Kristal II: Padatan Ion, Kovalen, dan Molekuler}

\begin{solution}{exo:b1:crystals-ionic-covalent:1}
Anion: $8 \times \frac18 + 6 \times \frac12 = 4$; kation: $12 \times
\frac14 + 1 = 4$. Setiap ion mempunyai 6 tetangga ion jenis lain (6:6).
Empat kation dan empat anion: \ce{NaCl}.
\end{solution}

\begin{solution}{exo:b1:crystals-ionic-covalent:2}
Setengah diagonal ruang: $a\sqrt3/2 = 412.3 \times 0.866 = \qty{357.1}{pm}$.
Satu \ce{Cs+} dan $8 \times \frac18 = 1$ \ce{Cl-}: satu satuan rumus.
\end{solution}

\begin{solution}{exo:b1:crystals-ionic-covalent:3}
$x = 102/181 = 0.56$, di antara 0.414 dan 0.732: koordinasi oktahedral
(6), \omterm{def:b1:crystals-ionic-covalent:structure-type}{tipe garam batu}.
\end{solution}

\begin{solution}{exo:b1:crystals-ionic-covalent:4}
Logam: tembaga. Ion: natrium klorida, kalsium fluorida. Kovalen: intan,
kuarsa. Molekuler: diiodin, es.
\end{solution}

\begin{solution}{exo:b1:crystals-ionic-covalent:5}
Sepanjang rusuk, anion--kation--anion bersentuhan: $a = 2(r_+ + r_-)$.
Sepanjang diagonal muka terletak dua anion pada jarak $a\sqrt2/2$;
keduanya tidak boleh bertumpang tindih: $a\sqrt2/2 \ge 2r_-$, sehingga
$(r_+ + r_-)\sqrt2 \ge 2r_-$, yaitu $r_+/r_- \ge \sqrt2 - 1$.
\end{solution}

\begin{solution}{exo:b1:crystals-ionic-covalent:6}
$\rho = 4 \times 58.44/(6.022 \times 10^{23} \times (5.641 \times
10^{-8})^3) = \qty{2.16}{g/cm^3}$, dibandingkan dengan \qty{2.17}{g/cm^3}
terukur.
\end{solution}

\begin{solution}{exo:b1:crystals-ionic-covalent:7}
$\rho = 168.36/(6.022 \times 10^{23} \times (4.123 \times 10^{-8})^3) =
\qty{3.99}{g/cm^3}$, sama dengan nilai terukurnya.
\end{solution}

\begin{solution}{exo:b1:crystals-ionic-covalent:8}
\ce{Zn-S} $= a\sqrt3/4 = \qty{234.2}{pm}$;
$\rho = 4 \times 97.44/(6.022 \times 10^{23} \times (5.409 \times
10^{-8})^3) = \qty{4.09}{g/cm^3}$, mendekati nilai terukur
\qty{4.04}{g/cm^3} untuk sfalerit alam.
\end{solution}

\begin{solution}{exo:b1:crystals-ionic-covalent:9}
\ce{C-C} $= 245.6/\sqrt3 = \qty{141.8}{pm}$; antarlapisan $669.6/2 =
\qty{334.8}{pm}$. Volume sel $\frac{\sqrt3}{2}a^2c = \qty{3.498e-23}{cm^3}$,
$\rho = 4 \times 12.01/(6.022 \times 10^{23} \times 3.498 \times 10^{-23})
= \qty{2.28}{g/cm^3}$, dibandingkan dengan \qty{3.52}{g/cm^3} untuk
intan: lapisan-lapisannya berjauhan dan hanya diikat oleh gaya London.
\end{solution}

\begin{solution}{exo:b1:crystals-ionic-covalent:10}
Lihat bukti \cref{prop:b1:crystals-ionic-covalent:diamond}:
$C = \pi\sqrt3/16 = 0.34$. Kekerasan berasal dari kekuatan dan kekakuan
tiga dimensi ikatan-ikatan kovalennya, bukan dari kerapatan susunannya:
setiap atom hanya mempunyai empat tetangga, pada sudut-sudut yang tetap,
sehingga tersisa banyak ruang kosong.
\end{solution}

\begin{solution}{exo:b1:crystals-ionic-covalent:11}
\ce{Si-Si} $= a\sqrt3/4 = \qty{235.2}{pm}$, \omterm{def:b1:periodicity:radius}{jari-jari kovalen}
\qty{117.6}{pm}; $\rho = 8 \times 28.09/(6.022 \times 10^{23} \times
(5.431 \times 10^{-8})^3) = \qty{2.33}{g/cm^3}$, sama dengan nilai
terukurnya.
\end{solution}

\begin{solution}{exo:b1:crystals-ionic-covalent:12}
$x = 152/181 = 0.84 > 0.732$: aturan itu meramalkan \omterm{def:b1:crystals-ionic-covalent:structure-type}{tipe sesium klorida}.
Pada struktur garam batu, \ce{Rb-Cl} $= a/2 = \qty{329.1}{pm}$, mendekati
$152 + 181 = \qty{333}{pm}$: struktur yang teramati konsisten. Aturan itu
memperlakukan ion sebagai bola keras dan mengabaikan \omterm{def:b1:periodicity:polarisability}{polarisabilitasnya}
serta selisih energi yang kecil antara kedua struktur (rubidium klorida
mengambil \omterm{def:b1:crystals-ionic-covalent:structure-type}{tipe sesium klorida} di bawah tekanan); aturan itu pedoman,
bukan hukum.
\end{solution}

\begin{solution}{pb:b1:crystals-ionic-covalent:1}
\textbf{1.} Di kedelapan sudut dan keenam pusat muka: $8 \times \frac18 + 6
\times \frac12 = 4$ ion.
\textbf{2.} 8 \omterm{def:b1:crystals-metals:site}{situs tetrahedral}, semuanya terisi: 8 \ce{F-} untuk 4
\ce{Ca^{2+}}, rasio 1:2, \ce{CaF2}.
\textbf{3.} \ce{F-}: 4 \ce{Ca^{2+}} di sudut-sudut sebuah tetrahedron.
\ce{Ca^{2+}}: 8 \ce{F-} di sudut-sudut sebuah kubus.
\textbf{4.} $a\sqrt3/4 = 546.3 \times 0.4330 = \qty{236.6}{pm}$.
\textbf{5.} $112 + 131 = \qty{243}{pm}$: selisihnya kurang dari 3\,\%.
\textbf{6.} Ion-ion \ce{F-} membentuk susunan kubus sederhana berusuk
$a/2 =
\qty{273.2}{pm}$, sedikit lebih besar daripada
$2 \times 131 = \qty{262}{pm}$: ion-ion itu hampir bersentuhan.
\textbf{7.} $x = 112/131 = 0.85$.
\textbf{8.} $x > 0.732$: koordinasi 8.
\textbf{9.} \omterm{def:b1:crystals-ionic-covalent:structure-type}{Tipe sesium klorida} mempunyai kation sama banyak dengan
anion; dengan anion dua kali lebih banyak, hanya separuh kubus-kubus
anion yang dapat memuat kation.
\textbf{10.} Kedelapan \ce{F-} sel itu membentuk kubus yang terdiri atas
8 kubus kecil berusuk $a/2$; ion-ion \ce{Ca^{2+}} menempati pusat 4 di
antaranya, secara berselang-seling, seperti petak-petak hitam papan
catur tiga dimensi.
\textbf{11.} $C = \frac43\pi(4 \times 112^3 + 8 \times 131^3)/546.3^3 = 0.61$.
\textbf{12.} \ce{O^{2-}} pada \omterm{def:b1:crystals-metals:lattice}{kisi} FCC (sudut-sudut dan pusat-pusat
muka), \ce{Li+} di kedelapan \omterm{def:b1:crystals-metals:site}{situs tetrahedral}.
\textbf{13.} $a\sqrt3/4 = \qty{200.0}{pm}$; $59 + 142 = \qty{201}{pm}$.
\textbf{14.} $\rho = 4 \times 29.88/(6.022 \times 10^{23} \times (4.619
\times 10^{-8})^3) = \qty{2.01}{g/cm^3}$, sama dengan nilai terukurnya.
\textbf{15.} Sfalerit: intan adalah sfalerit dengan karbon pada kedua
jenis posisi.
\textbf{16.} $a\sqrt3/4 = \qty{154.5}{pm}$; $\rho = 8 \times
12.01/(6.022 \times 10^{23} \times (3.567 \times 10^{-8})^3) =
\qty{3.52}{g/cm^3}$.
\textbf{17.} Pada intan hanya separuh \omterm{def:b1:crystals-metals:site}{situs tetrahedral} yang terisi, dan
oleh atom-atom yang sama besar dengan atom-atom \omterm{def:b1:crystals-metals:lattice}{kisi}, yang hanya
bersentuhan sepanjang ikatan: $C = 0.34$. Pada fluorit semua situs
terisi dan kation-kation kecil mengisi celah di antara anion-anion
besar: $C = 0.61$.
\textbf{18.} $40.08 + 2 \times 19.00 = \qty{78.08}{g/mol}$.
\textbf{19.} $4 \times 78.08/(6.022 \times 10^{23}) = \qty{5.186e-22}{g}$.
\textbf{20.} $(5.463 \times 10^{-8})^3 = \qty{1.630e-22}{cm^3}$.
\textbf{21.} Satu \ce{F-} yang hilang per seratus sel mengurangi massa
sebesar $19.00/(100
\times 312.3)$, 0.06\,\%: tidak terlihat pada tiga
angka (dan satu \ce{Ca^{2+}} juga harus hilang, kalau tidak \omterm{def:b1:crystals-metals:crystal}{kristal} itu
tidak netral).
\textbf{22.} $\rho = 5.186 \times 10^{-22}/1.630 \times 10^{-22} =
\textbf{\qty{3.18}{g/cm^3}}$, tepat sama dengan massa jenis terukurnya:
model sel, yang dibangun dari satu pengukuran $a$ dengan sinar-X,
menjelaskan massa mineral itu.
\end{solution}
